\section{基数上的序关系和运算}

\begin{proposition}{}{}
	设$\mathfrak{a},\mathfrak{b}$是基数,则$\mathfrak{a}\leq\mathfrak{b}$ $\iff$存在基数$\mathfrak{c}$使得$\mathfrak{b}+\mathfrak{c}=\mathfrak{a}$.
\end{proposition}

\begin{proposition}{}{}
	设$\left(\mathfrak{a}_{i}\right)_{i\in I},\left(\mathfrak{b}_{i}\right)_{i\in I}$是基数族.若$\forall i\in I\left(\mathfrak{a}_{i}\leq\mathfrak{b}_{i}\right)$,则$\sum\limits_{i\in I}\mathfrak{a}_{i}\leq\sum\limits_{i\in I}\mathfrak{b}_{i},\prod\limits_{i\in I}\mathfrak{a}_{i}\leq\prod\limits_{i\in I}\mathfrak{b}_{i}$.
\end{proposition}

\begin{corollary}{}{}
	设$\left(\mathfrak{a}_{i}\right)_{i\in I}$是基数族,$J\subseteq I$.
	\begin{enumerate}
		\item $\sum\limits_{i\in J}\mathfrak{a}_{i}\leq\sum\limits_{i\in I}\mathfrak{a}_{i}$.
		\item 若$\forall i\in I\setminus J\left(\mathfrak{a}_{i}\neq 0\right)$,则$\prod\limits_{i\in J}\mathfrak{a}_{i}\leq\prod\limits_{i\in I}\mathfrak{a}_{i}$.
	\end{enumerate}
\end{corollary}

\begin{corollary}{}{}
	设$\mathfrak{a},\mathfrak{b},\mathfrak{c},\mathfrak{d}$是基数.若$\mathfrak{a}\leq\mathfrak{b},\mathfrak{c}\leq\mathfrak{d},\mathfrak{b}>0$,则$\mathfrak{a}^{\mathfrak{b}}\leq\mathfrak{c}^{\mathfrak{d}}$.
\end{corollary}

\begin{theorem}{Cantor定理}{}
	对任意基数$\mathfrak{a}$有$2^{\mathfrak{a}}>\mathfrak{a}$.
\end{theorem}

\begin{corollary}{Burali-Forti悖论}{}
	不存在由所有的基数组成的集合.
\end{corollary}


























